\section{From archimedean resonance to a product window}
\label{sec:resonance-carrier}

Let \(q\) be an odd prime and identify each nonzero residue with its
unique signed representative in \((-q/2,q/2)\).  For
\(1\le H\le(q-1)/2\), define
\begin{equation}
 E_H
 =
 \{x\in\mathbb F_q^\times:|\widetilde x|\le H\}
 =
 \{\pm1,\ldots,\pm H\}.
\label{eq:EH}
\end{equation}
Thus \(|E_H|=2H\).

For \(1\le N\le q\), put
\begin{equation}
 H_q(N)
 =
 \min\!\left\{\frac{q-1}{2},\left\lfloor\frac qN\right\rfloor\right\}.
\label{eq:HqN}
\end{equation}

\begin{lemma}[Exact resonance window]
\label{lem:exact-resonance-window}
Let \(1\le N\le q\) and \(r,\omega\in\mathbb F_q^\times\).  Then
\[
 N\left\|\frac{\widetilde{r\omega}}q\right\|\le1
 \quad\Longleftrightarrow\quad
 r\omega\in E_{H_q(N)}.
\]
\end{lemma}

\begin{proof}
For \(N\ge2\), the left side is equivalent to
\[
 |\widetilde{r\omega}|\le q/N.
\]
The absolute residue is a positive integer, so this is equivalent to
\(|\widetilde{r\omega}|\le\lfloor q/N\rfloor=H_q(N)\).
For \(N=1\), both sides hold for every nonzero residue because
\(H_q(1)=(q-1)/2\).
\end{proof}

\begin{remark}[The range \(N>q\)]
If \(N>q\), no nonzero residue can be resonant because
\(|\widetilde{r\omega}|\ge1>q/N\).  Thus
\cref{lem:exact-resonance-window} also recovers the conductor--length
bound from TPC-94.
\end{remark}

\begin{definition}[Resonance incidence operator]
\label{def:KH}
On \(G=\mathbb F_q^\times\), set
\begin{equation}
 (K_Hb)(\omega)
 =
 \sum_{r\in G}b(r)\one_{E_H}(r\omega).
\label{eq:KH}
\end{equation}
For \(a,b:G\to\mathbb C\), define
\begin{equation}
 \mathcal I_H(a,b)
 =
 \sum_{\omega\in G}a(\omega)(K_Hb)(\omega).
\label{eq:incidence-form}
\end{equation}
\end{definition}

\begin{proposition}[Biregularity and total incidence]
\label{prop:biregular}
The bipartite graph
\[
 r\sim\omega
 \quad\Longleftrightarrow\quad
 r\omega\in E_H
\]
is \(2H\)-regular on both copies of \(G\).  In particular, for every
\(S\subseteq G\),
\[
 \sum_{\omega\in G}
 \#\{r\in S:r\omega\in E_H\}
 =2H|S|.
\]
\end{proposition}

\begin{proof}
For fixed \(r\), multiplication by \(r\) is a bijection and exactly
\(2H\) values of \(\omega\) send \(r\omega\) into \(E_H\).  The same
argument with the roles reversed proves both statements.
\end{proof}
